% SGA 3，第 VII_B 篇，英文重建稿——第 1 轮。
% 范围：权威底本分篇第 90 页，第 5.2.1 节开头及归约命题全文，止于第 5.2.1.A 节之前。
% 合并阅读本页码： 603.
% 再版印刷页码： 592.
% 本范围内原版页边定位： 558.
% 权威底本： Polo--Gille Exp7B-13oct24.pdf,
% SHA-256 D303690E0195654BBE7C3DAA5D89EA23E24CD709A3575217234AC19AE8836694.
% 曾参考 Jacob Reinhold 的 CC BY 4.0 Markdown 作为比对谱系。
% 本组件没有图表；全部数学均为原生 LaTeX。

\subsubsection*{5.2.1}
\phantomsection
\label{sga3:VIIB:5.2.1}

令 \(A\) 为 \(G\) 的仿射代数，\(I_A\) 为其增广理想，并置
\[
  I=B\cap I_A.
\]
由 \SGARefLink{sga3:VIIB:2.4}{2.4}，有
\[
  H=\operatorname{Spf}\bigl(A/\overline{AI}\bigr)
\]
及
\[
  B
  =
  \left\{
    x\in A\ \middle|\
    \Delta(x)-x\widehat\otimes1
    \in A\widehat\otimes\overline{AI}
  \right\}.
\]
\phantomsection
\label{sga3:VIIB:5.2.1:tensor-order-source-note}
\noindent\textit{底本说明。}\enspace
法文 PDF 此处印作 \(\Delta(x)-1\widehat\otimes x\)。\SGARefLink{sga3:VIIB:5.0}{5.0} 和 \SGARefLink{sga3:VIIB:5.1:proposition}{命题~5.1} 中确定的右商约定，连同定义条件 \((\operatorname{id}_A\widehat\otimes\pi)\Delta(x)
=x\widehat\otimes1\)，要求使用 \(\Delta(x)-x\widehat\otimes1\)；上文采用了后者。

置
\[
  \omega=I/\overline{I^2}.
\]
对每个 \(r\)，用 \(I_r\) 表示 \(I\) 中由满足
\[
  x^{p^r}=0,
\]
的元素 \(x\) 组成的闭子理想，用 \(\omega_r\) 表示 \(I_r\) 在 \(\omega\) 中的典范像。我们将证明如下命题。

\medskip
\noindent\textbf{命题。}
\phantomsection
\label{sga3:VIIB:5.2.1:section-proposition}
{\itshape 若投影 \(I\to I/\overline{I^2}\) 有连续截面
\[
  \sigma:\omega\longrightarrow I
\]
使得对每个 \(r\)，有
\[
  \sigma(\omega_r)\subset I_r
\]
则 \(B\) 同构于
\[
  k[[\omega]]
  \big/
  \bigl((x^{p^r})\bigr)_{x\in\omega_r}.
\]}

事实上，这样的截面延拓为连续态射
\[
  k[[\omega]]\longrightarrow B,
\]
它经由
\[
  B'
  =
  k[[\omega]]
  \big/
  \bigl((x^{p^r})\bigr)_{x\in\omega_r}.
\]
分解。在 \SGARefLink{sga3:VIIB:5.2.2}{5.2.2}--\SGARefLink{sga3:VIIB:5.2.5}{5.2.5} 中，我们证明所得态射
\[
  \varphi:
  B'
  =
  k[[\omega]]
  \big/
  \bigl((x^{p^r})\bigr)_{x\in\omega_r}
  \longrightarrow B
\]
为同构。
